%% Schéma cinématique d'un engrenage à deux roues (transmission de mouvement).
%% Roue 1 menante (petite, rouge) et roue 2 menée (grande, bleue) : cercles primitifs
%% tangents, axes parallèles, chaque roue montée en pivot sur le même bâti (noir).
%% Le rapport de transmission vaut R = w2/w1 = Z1/Z2 ; les deux roues tournent en sens opposés.
\documentclass[tikz,border=4pt]{standalone}
\usepackage{liaisons}
\usetikzlibrary{backgrounds}
\begin{document}
\begin{tikzpicture}[x=1cm,y=1cm,
    roue/.style={line width=1.4pt, line cap=round, line join=round},
    arbre/.style={line width=1.1pt, line cap=round},
    fleche/.style={-{Stealth[length=5pt]}, line width=0.8pt}]
\useasboundingbox (-3.3,-2.6) rectangle (3.5,3.7);   % boîte fixe (animation : cm → SVG)
% --- Macro locale : segment de bâti vertical hachuré (hachures vers la gauche) ---
\newcommand\bativerti[4]{% #1 = x, #2 = y bas, #3 = y haut, #4 = nombre de hachures
  \draw[line width=0.8pt] (#1,#2) -- (#1,#3);
  \foreach \i in {1,...,#4} {\pgfmathsetmacro\yy{#2+(#3-#2)*\i/#4}%
    \draw[line width=0.5pt] (#1,\yy) -- (#1-0.16,\yy-0.16);}}
% --- Géométrie : centres alignés verticalement, cercles primitifs tangents en I ---
\coordinate (O1) at (0,2.2);   % roue 1 (menante), rayon 0,8
\coordinate (O2) at (0,0);     % roue 2 (menée),   rayon 1,4
\coordinate (I)  at (0,1.4);   % point de contact (tangence des cercles primitifs)
% --- Bâti : montant vertical à gauche + bras vers chaque axe ---
\bativerti{-2.5}{-0.5}{2.7}{17}
\begin{scope}[on background layer]
  \draw[arbre] (-2.5,2.2) -- (O1);
  \draw[arbre] (-2.5,0)   -- (O2);
\end{scope}
% --- Roues : cercles primitifs + un rayon (liaison roue/arbre) ---
\draw[roue, solideA] (O1) circle (0.8);
\draw[roue, solideB] (O2) circle (1.4);
\draw[roue, solideA] (O1) -- ++(35:0.8);
\draw[roue, solideB] (O2) -- ++(-35:1.4);
% --- Dentures suggérées de part et d'autre du point de contact ---
\foreach \a in {250,262,...,290} {\draw[line width=0.7pt, solideA] ([shift=(\a:0.72)]O1) -- ([shift=(\a:0.9)]O1);}
\foreach \a in {73,80,...,108} {\draw[line width=0.7pt, solideB] ([shift=(\a:1.32)]O2) -- ([shift=(\a:1.5)]O2);}
\fill (I) circle (1.3pt);
\node[font=\small, anchor=west, inner sep=3pt] at (I) {I};
% --- Liaisons pivot d'axe z (vue « bout ») ---
\pic[s1=solideA, s2=black, liens=false] at (O1) {pivot bout};
\pic[s1=solideB, s2=black, liens=false] at (O2) {pivot bout};
% --- Sens de rotation : opposés ---
\draw[fleche, solideA!60] ([shift=(-20:1.15)]O1) arc (-20:35:1.15)
  node[pos=1, anchor=west, inner sep=2pt, font=\small, text=solideA] {$\omega_1$};
\draw[fleche, solideB!60] ([shift=(40:1.75)]O2) arc (40:-25:1.75)
  node[pos=1, anchor=west, inner sep=2pt, font=\small, text=solideB] {$\omega_2$};
% --- Étiquettes ---
\node[font=\small, text=solideA, anchor=south] at (0,3.15) {\textbf{1} menante ($Z_1$)};
\node[font=\small, text=solideB, anchor=north] at (0,-1.6) {\textbf{2} menée ($Z_2$)};
\node[font=\small, anchor=east] at (-2.75,1.1) {\textbf{0}};
\repere{(2.25,-2.05)}{x}{y}{1}
\end{tikzpicture}
\end{document}
